\section{Residual context topology and sheaf-local Birkhoff}
\label{sec:residual-context-topology}

The global coequational theorem can be sharpened by separating two levels of
globality. The first is stability under residual restriction. The second is
descent along covers of residual contexts. This section constructs the
canonical residual-context topology and records the corresponding sheaf-local
refinement of the global Birkhoff theorem.

Throughout this section the output internal category \(\mathbb B\) is fixed.
The category of allowed input categories and input functors is denoted
\[
    \mathsf{Inp}.
\]
We work in full Grothendieck-topos generality by using generalized residual
contexts. Thus we do not assume that objects have enough global points. All
indexing categories used below are understood to be small relative to the fixed
ambient universe. If the collection of all generalized residual contexts is
large, one restricts to a chosen small full site of contexts, closed under the
pullbacks and restrictions used below, or to a small set of representatives.

\subsection{Generalized residual context shape}
\label{subsec:generalized-residual-context-shape}

\begin{definition}[Generalized residual context category]
\label{def:generalized-residual-context-category}
    The \textbf{generalized residual context category}
    \[
        \mathsf{RCtx}(\mathsf{Inp})
    \]
    is defined as follows.

    An object is a generalized pointed input context
    \[
        (\mathbb A,U,\nu),
        \qquad
        \mathbb A\in\mathsf{Inp},
        \qquad
        \nu:U\to A_0.
    \]
    Thus \(U\) is a parameter object and \(\nu\) is a \(U\)-indexed family of
    input objects.

    A morphism
    \[
        (F,h,r):
        (\mathbb A',U',\nu')
        \longrightarrow
        (\mathbb A,U,\nu)
    \]
    consists of an input functor
    \[
        F:\mathbb A'\to\mathbb A,
    \]
    a parameter map
    \[
        h:U'\to U,
    \]
    and a \(U'\)-indexed arrow of \(\mathbb A\)
    \[
        r:U'\to A_1
    \]
    satisfying
    \[
        s_A r=F_0\nu',
        \qquad
        t_A r=\nu h.
    \]
    In generalized-element notation, \(r\) is a family of arrows
    \[
        r_{u'}:F_0(\nu'(u'))\to \nu(h(u')).
    \]
\end{definition}

The typing of \(r\) is displayed by
\[
\begin{tikzcd}[column sep=large, row sep=large]
    &
    A_1
        \arrow[d, "{(s_A,t_A)}"]
    \\
    U'
        \arrow[ur, "r"]
        \arrow[r, "{(F_0\nu',\,\nu h)}"']
    &
    A_0\times A_0 .
\end{tikzcd}
\]

The identity on
\[
    (\mathbb A,U,\nu)
\]
is
\[
    (1_{\mathbb A},1_U,e_A\nu).
\]
If
\[
    (G,k,s):
    (\mathbb A'',U'',\nu'')
    \to
    (\mathbb A',U',\nu')
\]
and
\[
    (F,h,r):
    (\mathbb A',U',\nu')
    \to
    (\mathbb A,U,\nu),
\]
their composite is
\[
    (F,h,r)\circ(G,k,s)
    :=
    \bigl(FG,\;hk,\;m_A(F_1s,rk)\bigr).
\]
Here
\[
    rk:=r\circ k:U''\to A_1,
\]
and
\[
    F_1s:U''\to A_1
\]
is the image under \(F\) of the \(U''\)-indexed arrow \(s\). In
generalized-element notation, the composite arrow is
\[
    F_0G_0(\nu''(u''))
        \xrightarrow{F_1(s_{u''})}
    F_0(\nu'(k u''))
        \xrightarrow{r_{k u''}}
    \nu(hk u'').
\]
Thus
\[
    m_A(F_1s,rk)=rk\circ F_1s.
\]

The composite-arrow typing is displayed by
\[
\begin{tikzcd}[column sep=large, row sep=large]
    F_0G_0\nu''
        \arrow[r, "{F_1s}"]
        \arrow[dr, swap, "{m_A(F_1s,rk)}"]
    &
    F_0\nu'k
        \arrow[d, "rk"]
    \\
    &
    \nu hk .
\end{tikzcd}
\]

\begin{lemma}[Residual context category]
\label{lem:residual-context-category}
    The data above define a category
    \[
        \mathsf{RCtx}(\mathsf{Inp}).
    \]
\end{lemma}

\begin{proof}
    \leavevmode
    \begin{enumerate}
        \item \textbf{Check identities.}
        Let
        \[
            (F,h,r):
            (\mathbb A',U',\nu')
            \to
            (\mathbb A,U,\nu)
        \]
        be a residual context morphism. Composing on the source side with
        \[
            (1_{\mathbb A'},1_{U'},e_{A'}\nu')
        \]
        gives the arrow component
        \[
            m_A(F_1(e_{A'}\nu'),r)
            =
            m_A(e_A F_0\nu',r)
            =
            r.
        \]
        Composing on the target side with
        \[
            (1_{\mathbb A},1_U,e_A\nu)
        \]
        gives the arrow component
        \[
            m_A(r,e_A\nu h)=r.
        \]
        These are precisely the identity laws in \(\mathbb A\).

        \item \textbf{Check associativity.}
        Given three composable residual context morphisms, the parameter-map
        part is associative because composition in \(\mathcal E\) is
        associative, and the input-functor part is associative because
        composition in \(\mathsf{Inp}\) is associative.

        For the arrow component, associativity reduces to associativity of
        composition in \(\mathbb A\). If the three displayed arrow components
        are
        \[
            t,\qquad s,\qquad r,
        \]
        then the two possible composites have arrow components
        \[
            (rk)\circ F_1(s\ell)\circ F_1G_1(t)
        \]
        and
        \[
            rk\circ F_1(s\ell\circ G_1(t)).
        \]
        These agree by functoriality of \(F\) and associativity in
        \(\mathbb A\).

        \item \textbf{Display the associativity calculation.}
        The relevant pasting diagram is
        \[
        \begin{tikzcd}[column sep=large, row sep=large]
            F_0G_0H_0\nu'''
                \arrow[r, "{F_1G_1(t)}"]
                \arrow[rr, bend left=14, "{F_1(G_1(t)\circ s\ell)}"]
                \arrow[rrr, bend left=26]
            &
            F_0G_0\nu''\ell
                \arrow[r, "{F_1(s\ell)}"]
            &
            F_0\nu'k\ell
                \arrow[r, "{rk\ell}"]
            &
            \nu hk\ell .
        \end{tikzcd}
        \]
        Pasting either the first two arrows and then the third, or the last two
        arrows and then the first, gives the same composite.
    \end{enumerate}
\end{proof}

For a generalized context
\[
    (\mathbb A,U,\nu),
\]
write
\[
    \mathbb A/\nu
\]
for the residual slice internal category in the slice
\[
    \mathcal E/U.
\]
Its object object is
\[
    U\times_{\nu,A_0,t_A}A_1\to U,
\]
and its arrow object is
\[
    U\times_{\nu,A_0,t_A\pi_2}T_A\to U,
    \qquad
    T_A:=A_1\times_{t_A,A_0,s_A}A_1.
\]
A generalized object is a pair
\[
    (u,a:x\to\nu(u)),
\]
and a generalized residual arrow is a triangle
\[
\begin{tikzcd}[column sep=large, row sep=large]
    y
        \arrow[r, "c"]
        \arrow[rr, "a\circ c", bend left=15]
    &
    x
        \arrow[r, "a"']
    &
    \nu(u).
\end{tikzcd}
\]

A residual context morphism
\[
    (F,h,r):
    (\mathbb A',U',\nu')
    \to
    (\mathbb A,U,\nu)
\]
induces an internal functor in \(\mathcal E/U'\)
\[
    (F,h,r)_{/}:
    \mathbb A'/\nu'
    \longrightarrow
    h^*(\mathbb A/\nu).
\]
On residual objects it is given by
\[
    a':x'\to\nu'(u')
    \quad\longmapsto\quad
    r_{u'}\circ F_1(a'):F_0(x')\to\nu(h u').
\]
On residual arrows it sends the triangle
\[
\begin{tikzcd}[column sep=large, row sep=large]
    y'
        \arrow[r, "c'"]
        \arrow[rr, "a'\circ c'", bend left=15]
    &
    x'
        \arrow[r, "a'"']
    &
    \nu'(u')
\end{tikzcd}
\]
to the triangle
\[
\begin{tikzcd}[column sep=large, row sep=large]
    F_0(y')
        \arrow[r, "F_1(c')"]
        \arrow[rr, "{r_{u'}\circ F_1(a'\circ c')}"', bend right=15]
    &
    F_0(x')
        \arrow[r, "{r_{u'}\circ F_1(a')}"']
    &
    \nu(h u').
\end{tikzcd}
\]

\begin{lemma}[Residual slice functoriality]
\label{lem:residual-slice-functoriality}
    The assignment
    \[
        (\mathbb A,U,\nu)\longmapsto \mathbb A/\nu
    \]
    and
    \[
        (F,h,r)\longmapsto (F,h,r)_{/}
    \]
    is functorial in the indexed sense: identities induce identity residual
    slice functors, and composites induce the corresponding composite residual
    slice functors.
\end{lemma}

\begin{proof}
    \leavevmode
    \begin{enumerate}
        \item \textbf{Check identities.}
        The identity residual context morphism
        \[
            (1_{\mathbb A},1_U,e_A\nu)
        \]
        sends a residual object
        \[
            a:x\to\nu(u)
        \]
        to
        \[
            e_A\nu(u)\circ a=a.
        \]
        It also sends each residual triangle to itself. Hence it induces the
        identity functor on
        \[
            \mathbb A/\nu.
        \]

        \item \textbf{Check composition on residual objects.}
        Let
        \[
            (G,k,s):
            (\mathbb A'',U'',\nu'')
            \to
            (\mathbb A',U',\nu')
        \]
        and
        \[
            (F,h,r):
            (\mathbb A',U',\nu')
            \to
            (\mathbb A,U,\nu)
        \]
        be composable. For a residual object
        \[
            a'':x''\to\nu''(u''),
        \]
        the two-step residual restriction sends \(a''\) to
        \[
            r_{k u''}\circ F_1(s_{u''})\circ F_1G_1(a'').
        \]
        The composite residual context morphism has arrow component
        \[
            r_{k u''}\circ F_1(s_{u''}),
        \]
        and therefore sends \(a''\) to the same arrow
        \[
            \bigl(r_{k u''}\circ F_1(s_{u''})\bigr)
            \circ F_1G_1(a'').
        \]

        \item \textbf{Check composition on residual arrows.}
        The same calculation applied to a residual triangle
        \[
        \begin{tikzcd}[column sep=large, row sep=large]
            y''
                \arrow[r, "c''"]
                \arrow[rr, "a''\circ c''", bend left=15]
            &
            x''
                \arrow[r, "a''"']
            &
            \nu''(u'')
        \end{tikzcd}
        \]
        shows that both composite functors send it to the triangle
        \[
        \begin{tikzcd}[column sep=large, row sep=large]
            F_0G_0(y'')
                \arrow[r, "{F_1G_1(c'')}"]
                \arrow[rr,
                    "{r_{k u''}\circ F_1(s_{u''})\circ F_1G_1(a''\circ c'')}"',
                    bend right=15]
            &
            F_0G_0(x'')
                \arrow[r,
                    "{r_{k u''}\circ F_1(s_{u''})\circ F_1G_1(a'')}"']
            &
            \nu(hk u'').
        \end{tikzcd}
        \]
        Therefore the residual slice assignment is functorial.
    \end{enumerate}
\end{proof}

The two restriction operations used earlier are special cases of residual
context restriction.

\begin{lemma}[Derivatives and input reindexing as residual restrictions]
\label{lem:derivatives-and-reindexing-as-residual-restrictions}
    Let
    \[
        H:U\to\mathbf{Beh}_{\mathbb A,\mathbb B,0}
    \]
    be a \(U\)-indexed residual behaviour over
    \[
        \nu:U\to A_0,
    \]
    so that
    \[
        p_{\mathbf{Beh}}H=\nu.
    \]

    \begin{enumerate}
        \item If
        \[
            a:U\to A_1
        \]
        is a \(U\)-indexed arrow of \(\mathbb A\) with
        \[
            s_Aa=\mu,
            \qquad
            t_Aa=\nu,
        \]
        then restriction along
        \[
            (1_{\mathbb A},1_U,a):
            (\mathbb A,U,\mu)\to(\mathbb A,U,\nu)
        \]
        is the \(U\)-indexed derivative
        \[
            H\longmapsto\partial_aH.
        \]

        \item If
        \[
            F:\mathbb A'\to\mathbb A
        \]
        is an input functor and
        \[
            \nu':U'\to A'_0
        \]
        is a generalized input object, then restriction along
        \[
            (F,1_{U'},e_AF_0\nu'):
            (\mathbb A',U',\nu')
            \to
            (\mathbb A,U',F_0\nu')
        \]
        is the input reindexing operation
        \[
            H\longmapsto F^\sharp(\nu',H).
        \]
    \end{enumerate}
\end{lemma}

\begin{proof}
    \leavevmode
    \begin{enumerate}
        \item \textbf{Derivatives.}
        The residual slice functor associated to
        \[
            (1_{\mathbb A},1_U,a)
        \]
        sends
        \[
            r:x\to\mu(u)
        \]
        to
        \[
            a_u\circ r:x\to\nu(u).
        \]
        This is precisely residual transport by the arrow \(a_u\). Therefore
        restriction of \(H\) along
        \[
            (1_{\mathbb A},1_U,a)
        \]
        is
        \[
            H\circ a_*=\partial_aH.
        \]

        \item \textbf{Input reindexing.}
        The residual slice functor associated to
        \[
            (F,1_{U'},e_AF_0\nu')
        \]
        sends
        \[
            r':x'\to\nu'(u')
        \]
        to
        \[
            F_1(r'):F_0(x')\to F_0(\nu'(u')).
        \]
        This is exactly the residual slice functor
        \[
            F_{/\nu'}:
            \mathbb A'/\nu'
            \to
            \mathbb A/F_0\nu'.
        \]
        Therefore restriction of \(H\) is
        \[
            H\circ F_{/\nu'}=F^\sharp(\nu',H).
        \]

        \item \textbf{Display the two special restrictions.}
        The derivative case is displayed by
        \[
        \begin{tikzcd}[column sep=large, row sep=large]
            x
                \arrow[r, "r"]
                \arrow[rr, "a_u\circ r", bend left=15]
            &
            \mu(u)
                \arrow[r, "a_u"']
            &
            \nu(u),
        \end{tikzcd}
        \]
        while the input-reindexing case is displayed by
        \[
        \begin{tikzcd}[column sep=large, row sep=large]
            x'
                \arrow[r, "r'"]
                \arrow[d, maps to]
            &
            \nu'(u')
                \arrow[d, maps to]
            \\
            F_0(x')
                \arrow[r, "F_1(r')"']
            &
            F_0(\nu'(u')).
        \end{tikzcd}
        \]
    \end{enumerate}
\end{proof}

Every residual context morphism factors as an input reindexing morphism
followed by a derivative-restriction morphism:
\[
\begin{tikzcd}[column sep=huge, row sep=large]
    (\mathbb A',U',\nu')
        \arrow[r, "{(F,1_{U'},e_AF_0\nu')}"]
        \arrow[rr, bend left=18, "{(F,h,r)}"]
    &
    (\mathbb A,U',F_0\nu')
        \arrow[r, "{(1_{\mathbb A},h,r)}"]
    &
    (\mathbb A,U,\nu).
\end{tikzcd}
\]
Thus
\[
    (F,h,r)
    =
    (1_{\mathbb A},h,r)\circ(F,1_{U'},e_AF_0\nu').
\]
Contravariantly, restriction along \((F,h,r)\) is derivative restriction
followed by input reindexing.

\subsection{The residual behaviour presheaf}
\label{subsec:residual-behaviour-presheaf}

\begin{definition}[Residual behaviour presheaf]
\label{def:residual-behaviour-presheaf}
    The \textbf{residual behaviour presheaf} with output \(\mathbb B\) is
    \[
        \mathcal B_{\mathbb B}:
        \mathsf{RCtx}(\mathsf{Inp})^{op}\to\mathbf{Set}
    \]
    defined by
    \[
        \mathcal B_{\mathbb B}(\mathbb A,U,\nu)
        :=
        \mathcal E/A_0
        \bigl((U,\nu),(\mathbf{Beh}_{\mathbb A,\mathbb B,0},
        p_{\mathbf{Beh}})\bigr).
    \]
    Equivalently, an element is a \(U\)-indexed residual behaviour
    \[
        H:U\to\mathbf{Beh}_{\mathbb A,\mathbb B,0}
    \]
    satisfying
    \[
        p_{\mathbf{Beh}}H=\nu.
    \]
    This is equivalently a \(U\)-indexed internal functor
    \[
        \mathbb A/\nu\to U^*\mathbb B
    \]
    in the slice \(\mathcal E/U\).

    For a residual context morphism
    \[
        (F,h,r):
        (\mathbb A',U',\nu')
        \to
        (\mathbb A,U,\nu),
    \]
    the restriction map
    \[
        (F,h,r)^*:
        \mathcal B_{\mathbb B}(\mathbb A,U,\nu)
        \to
        \mathcal B_{\mathbb B}(\mathbb A',U',\nu')
    \]
    sends \(H\) to the \(U'\)-indexed residual behaviour obtained by
    precomposition with
    \[
        (F,h,r)_{/}.
    \]
    In generalized-element notation,
    \[
        \bigl((F,h,r)^*H\bigr)(u')
        =
        H(h u')\circ(F,r_{u'})_{/}.
    \]
\end{definition}

The restriction operation is summarized by
\[
\begin{tikzcd}[column sep=huge, row sep=large]
    \mathbb A'/\nu'
        \arrow[r, "{(F,h,r)_{/}}"]
        \arrow[dr, "{(F,h,r)^*H}"']
    &
    h^*(\mathbb A/\nu)
        \arrow[d, "h^*H"]
    \\
    &
    U'^*\mathbb B .
\end{tikzcd}
\]

\begin{proposition}[Global theories as residual-context subpresheaves]
\label{prop:global-theories-as-residual-context-subpresheaves}
    Let
    \[
        V=(V_{\mathbb A})_{\mathbb A\in\mathsf{Inp}}
    \]
    be a family assigning to each input category \(\mathbb A\) a local
    coequational theory
    \[
        V_{\mathbb A}\hookrightarrow
        \mathbf{Beh}^{\mathrm{can}}_{\mathbb A,\mathbb B}.
    \]
    Then \(V\) is a global coequational theory precisely when the assignment
    \[
        (\mathbb A,U,\nu)
        \longmapsto
        \mathcal V(\mathbb A,U,\nu)
    \]
    defined by
    \[
        \mathcal V(\mathbb A,U,\nu)
        :=
        \mathcal E/A_0
        \bigl((U,\nu),(V_{\mathbb A},p_{V_{\mathbb A}})\bigr)
    \]
    is stable under all residual-context restrictions. Equivalently, it defines
    a subpresheaf
    \[
        \mathcal V\hookrightarrow\mathcal B_{\mathbb B}
    \]
    on
    \[
        \mathsf{RCtx}(\mathsf{Inp}).
    \]
\end{proposition}

\begin{proof}
    \leavevmode
    \begin{enumerate}
        \item \textbf{From global coequational theories to subpresheaves.}
        Suppose \(V\) is a global coequational theory. Local coequationality of
        each
        \[
            V_{\mathbb A}
        \]
        says that \(V_{\mathbb A}\) is closed under canonical derivatives.
        Since derivative-closedness is an internal factorization condition in
        the relevant carrier slice, it is stable under arbitrary parameter
        pullback. Therefore \(V_{\mathbb A}\) is stable under restrictions of
        the form
        \[
            (1_{\mathbb A},h,r)^*.
        \]

        The global input-reindexing condition says that \(V\) is stable under
        restrictions of the form
        \[
            (F,1_{U'},e_AF_0\nu')^*.
        \]

        Every residual context morphism factors as
        \[
            (F,h,r)
            =
            (1_{\mathbb A},h,r)
            \circ
            (F,1_{U'},e_AF_0\nu').
        \]
        Hence the family \(V\) is stable under restriction along every residual
        context morphism. Thus it defines a subpresheaf of
        \[
            \mathcal B_{\mathbb B}.
        \]

        \item \textbf{From subpresheaves to global coequational theories.}
        Conversely, suppose the fibrewise assignment defines a subpresheaf of
        \[
            \mathcal B_{\mathbb B}.
        \]
        Stability under restrictions along
        \[
            (1_{\mathbb A},1_U,a)
        \]
        gives derivative closure of each local theory
        \[
            V_{\mathbb A}.
        \]
        Stability under restrictions along
        \[
            (F,1_{U'},e_AF_0\nu')
        \]
        gives the global input-reindexing condition. Therefore \(V\) is a
        global coequational theory.

        \item \textbf{Display the restriction square.}
        The defining stability condition is the factorization
        \[
        \begin{tikzcd}[column sep=large, row sep=large]
            \mathcal V(\mathbb A,U,\nu)
                \arrow[r, hook]
                \arrow[d, dashed, "{(F,h,r)^*}"']
            &
            \mathcal B_{\mathbb B}(\mathbb A,U,\nu)
                \arrow[d, "{(F,h,r)^*}"]
            \\
            \mathcal V(\mathbb A',U',\nu')
                \arrow[r, hook]
            &
            \mathcal B_{\mathbb B}(\mathbb A',U',\nu')
        \end{tikzcd}
        \]
        for every residual context morphism
        \[
            (F,h,r):
            (\mathbb A',U',\nu')
            \to
            (\mathbb A,U,\nu).
        \]
    \end{enumerate}
\end{proof}

\subsection{The canonical residual behaviour topology}
\label{subsec:canonical-residual-behaviour-topology}

The residual behaviour presheaf determines a canonical topology on the
residual context category. This topology is canonical relative to the presheaf
\[
    \mathcal B_{\mathbb B};
\]
it is not asserted to be the canonical subcanonical topology on
\[
    \mathsf{RCtx}(\mathsf{Inp}).
\]
It is the largest topology for which residual behaviours satisfy descent.

Let
\[
    X:\mathcal C^{op}\to\mathbf{Set}
\]
be a presheaf on a small category \(\mathcal C\). For a sieve
\[
    S\hookrightarrow y(c)
\]
on \(c\), write
\[
    \operatorname{Match}(S,X)
    :=
    \operatorname{Nat}(S,X)
\]
for the set of matching families of \(X\)-sections along \(S\). The canonical
restriction map is
\[
    X(c)\to\operatorname{Match}(S,X).
\]
A sieve \(S\) is called \textbf{\(X\)-effective} if this map is a bijection.
It is called \textbf{universally \(X\)-effective} if, for every morphism
\[
    f:d\to c,
\]
the pullback sieve
\[
    f^*S
\]
is \(X\)-effective on \(d\).

\begin{lemma}[Topology generated by a presheaf]
\label{lem:topology-generated-by-a-presheaf}
    Let
    \[
        X:\mathcal C^{op}\to\mathbf{Set}
    \]
    be a presheaf. The universally \(X\)-effective sieves form a Grothendieck
    topology
    \[
        J_X
    \]
    on \(\mathcal C\). Moreover, \(J_X\) is the largest Grothendieck topology
    on \(\mathcal C\) for which \(X\) is a sheaf.
\end{lemma}

\begin{proof}
    \leavevmode
    \begin{enumerate}
        \item \textbf{Maximal sieves cover.}
        For every object \(c\), the maximal sieve is the representable
        presheaf
        \[
            y(c).
        \]
        By the Yoneda lemma,
        \[
            \operatorname{Nat}(y(c),X)\cong X(c),
        \]
        so the maximal sieve is \(X\)-effective. The same argument applies
        after arbitrary pullback. Hence maximal sieves are universally
        \(X\)-effective.

        \item \textbf{Pullback stability.}
        Pullback stability is built into universal \(X\)-effectivity. If \(S\)
        is universally \(X\)-effective on \(c\) and
        \[
            f:d\to c
        \]
        is any morphism, then \(f^*S\) is universally \(X\)-effective on \(d\).

        \item \textbf{Upward closure.}
        Let
        \[
            S\subseteq T
        \]
        be sieves on \(c\), and suppose \(S\) is universally \(X\)-effective.
        Let
        \[
            (x_t)_{t\in T}
        \]
        be a matching family over \(T\). Restricting to \(S\), there is a
        unique
        \[
            x\in X(c)
        \]
        whose restrictions along arrows of \(S\) are the given \(x_s\).

        For an arrow
        \[
            t:d\to c
        \]
        in \(T\), compare
        \[
            X(t)(x)
            \qquad\text{and}\qquad
            x_t
        \]
        in \(X(d)\). These two elements agree after restriction along every
        arrow of the pullback sieve
        \[
            t^*S.
        \]
        Since \(t^*S\) is \(X\)-effective, they agree in \(X(d)\). Thus \(x\)
        amalgamates the matching family over \(T\). Uniqueness follows because
        \(S\subseteq T\). The same argument after arbitrary pullback shows that
        \(T\) is universally \(X\)-effective.

        The comparison is displayed by
        \[
        \begin{tikzcd}[column sep=large, row sep=large]
            t^*S
                \arrow[r, hook]
                \arrow[d, hook]
            &
            y(d)
                \arrow[d, "y(t)"]
            \\
            S
                \arrow[r, hook]
            &
            y(c).
        \end{tikzcd}
        \]

        \item \textbf{Transitivity.}
        Let \(S\) be universally \(X\)-effective on \(c\). Let \(R\) be a
        sieve on \(c\) such that
        \[
            f^*R
        \]
        is universally \(X\)-effective for every arrow
        \[
            f:d\to c
        \]
        in \(S\). We prove that \(R\) is universally \(X\)-effective.

        First take a matching family
        \[
            (y_g)_{g\in R}
        \]
        over \(R\). For each arrow
        \[
            f:d\to c
        \]
        in \(S\), the restriction of this matching family to \(f^*R\) has a
        unique amalgamation
        \[
            x_f\in X(d),
        \]
        since \(f^*R\) is \(X\)-effective. The family
        \[
            (x_f)_{f\in S}
        \]
        is a matching family over \(S\). Since \(S\) is \(X\)-effective, it has
        a unique amalgamation
        \[
            x\in X(c).
        \]

        For any arrow
        \[
            g:e\to c
        \]
        in \(R\), compare
        \[
            X(g)(x)
            \qquad\text{and}\qquad
            y_g
        \]
        in \(X(e)\). They agree after restriction along the pullback sieve
        \[
            g^*S.
        \]
        Indeed, for an arrow
        \[
            \ell:d\to e
        \]
        with
        \[
            g\ell\in S,
        \]
        the section \(x_{g\ell}\) is the unique amalgamation of the restricted
        matching family over
        \[
            (g\ell)^*R.
        \]
        Since \(g\in R\), the element \(y_g\) restricts to the same matching
        family. Hence
        \[
            X(\ell)X(g)(x)=X(\ell)(y_g).
        \]
        Because \(g^*S\) is \(X\)-effective, this implies
        \[
            X(g)(x)=y_g.
        \]
        Hence \(x\) amalgamates the matching family over \(R\).

        For uniqueness, let
        \[
            x,x'\in X(c)
        \]
        be two amalgamations over \(R\). For each
        \[
            f:d\to c
        \]
        in \(S\), the two restrictions
        \[
            X(f)(x),\quad X(f)(x')
        \]
        agree after restriction along
        \[
            f^*R.
        \]
        Since \(f^*R\) is \(X\)-effective, they agree in \(X(d)\). Since \(S\)
        is \(X\)-effective, \(x=x'\). Therefore \(R\) is \(X\)-effective.
        The same argument after arbitrary pullback proves universal
        \(X\)-effectivity.

        The transitivity pasting is summarized by
        \[
        \begin{tikzcd}[column sep=large, row sep=large]
            f^*R
                \arrow[r, hook]
                \arrow[d, hook]
            &
            y(d)
                \arrow[d, "y(f)"]
            \\
            R
                \arrow[r, hook]
            &
            y(c)
            &
            S
                \arrow[l, hook']
        \end{tikzcd}
        \]
        for each \(f\in S\).

        \item \textbf{Largest topology for which \(X\) is a sheaf.}
        By construction, every \(J_X\)-covering sieve is universally
        \(X\)-effective, so \(X\) is a sheaf for \(J_X\).

        Conversely, let \(J\) be any Grothendieck topology for which \(X\) is a
        sheaf. If
        \[
            S\in J(c),
        \]
        then \(S\) is \(X\)-effective. Since pullbacks of \(J\)-covering
        sieves are again \(J\)-covering, every pullback of \(S\) is
        \(X\)-effective. Thus every \(J\)-covering sieve is universally
        \(X\)-effective. Therefore
        \[
            J\subseteq J_X.
        \]
        Hence \(J_X\) is the largest topology for which \(X\) is a sheaf.
    \end{enumerate}
\end{proof}

\begin{definition}[Canonical residual behaviour topology]
\label{def:canonical-residual-behaviour-topology}
    The \textbf{canonical residual behaviour topology} associated to the output
    category \(\mathbb B\) is the Grothendieck topology
    \[
        J_{\mathbb B}^{\mathrm{res}}
        :=
        J_{\mathcal B_{\mathbb B}}
    \]
    on
    \[
        \mathsf{RCtx}(\mathsf{Inp})
    \]
    whose covering sieves are the universally
    \(\mathcal B_{\mathbb B}\)-effective sieves.
\end{definition}

Thus a sieve
\[
    S
\]
on a residual context
\[
    (\mathbb A,U,\nu)
\]
is \(J_{\mathbb B}^{\mathrm{res}}\)-covering precisely when, after arbitrary
residual-context pullback, residual behaviours over the target context are
uniquely recoverable from compatible families of their restrictions along the
arrows of the pulled-back sieve.

\begin{corollary}[Residual behaviours form a sheaf]
\label{cor:residual-behaviours-form-sheaf}
    The residual behaviour presheaf
    \[
        \mathcal B_{\mathbb B}
    \]
    is a sheaf on
    \[
        \bigl(\mathsf{RCtx}(\mathsf{Inp}),J_{\mathbb B}^{\mathrm{res}}\bigr).
    \]
    Moreover,
    \[
        J_{\mathbb B}^{\mathrm{res}}
    \]
    is the largest Grothendieck topology on
    \[
        \mathsf{RCtx}(\mathsf{Inp})
    \]
    with this property.
\end{corollary}

\begin{proof}
    This is Lemma~\ref{lem:topology-generated-by-a-presheaf} applied to
    \[
        X=\mathcal B_{\mathbb B}.
    \]
\end{proof}

\subsection{Sheaf-local global coequational theories}
\label{subsec:sheaf-local-global-coequational-theories}

\begin{definition}[Sheaf-local global coequational theory]
\label{def:sheaf-local-global-coequational-theory}
    A global coequational theory
    \[
        V=(V_{\mathbb A})_{\mathbb A\in\mathsf{Inp}}
    \]
    is \textbf{sheaf-local} if the corresponding residual-context subpresheaf
    \[
        \mathcal V\hookrightarrow\mathcal B_{\mathbb B}
    \]
    is a subsheaf for the canonical residual behaviour topology
    \[
        J_{\mathbb B}^{\mathrm{res}}.
    \]
\end{definition}

\begin{lemma}[Membership criterion for sheaf-local theories]
\label{lem:membership-criterion-sheaf-local-theories}
    Let
    \[
        V=(V_{\mathbb A})_{\mathbb A\in\mathsf{Inp}}
    \]
    be a global coequational theory. Then \(V\) is sheaf-local if and only if
    the following condition holds.

    For every residual context
    \[
        (\mathbb A,U,\nu),
    \]
    every residual behaviour
    \[
        H\in\mathcal B_{\mathbb B}(\mathbb A,U,\nu),
    \]
    and every covering sieve
    \[
        S\in J_{\mathbb B}^{\mathrm{res}}(\mathbb A,U,\nu),
    \]
    if
    \[
        (F,h,r)^*H
        \in
        \mathcal V(\mathbb A',U',\nu')
    \]
    for every arrow
    \[
        (F,h,r):
        (\mathbb A',U',\nu')
        \to
        (\mathbb A,U,\nu)
    \]
    in \(S\), then
    \[
        H\in\mathcal V(\mathbb A,U,\nu).
    \]
\end{lemma}

\begin{proof}
    \leavevmode
    \begin{enumerate}
        \item \textbf{Use the subpresheaf condition.}
        Since \(V\) is already a global coequational theory, the associated
        assignment
        \[
            \mathcal V\hookrightarrow\mathcal B_{\mathbb B}
        \]
        is a subpresheaf by
        Proposition~\ref{prop:global-theories-as-residual-context-subpresheaves}.
        Hence if
        \[
            H\in\mathcal V(\mathbb A,U,\nu),
        \]
        then every restriction
        \[
            (F,h,r)^*H
        \]
        lies in the corresponding fibre of \(\mathcal V\).

        \item \textbf{Unpack the subsheaf condition.}
        Since
        \[
            \mathcal B_{\mathbb B}
        \]
        is a sheaf for
        \[
            J_{\mathbb B}^{\mathrm{res}},
        \]
        a subpresheaf
        \[
            \mathcal V\hookrightarrow\mathcal B_{\mathbb B}
        \]
        is a subsheaf precisely when membership in \(\mathcal V\) is local on
        covering sieves. That is, whenever a section
        \[
            H\in\mathcal B_{\mathbb B}(\mathbb A,U,\nu)
        \]
        restricts to sections of \(\mathcal V\) along every arrow of a
        covering sieve, the section \(H\) itself lies in
        \[
            \mathcal V(\mathbb A,U,\nu).
        \]

        \item \textbf{Display the local membership test.}
        The membership test is the sheaf-local square
        \[
        \begin{tikzcd}[column sep=large, row sep=large]
            S
                \arrow[r, hook]
                \arrow[d, dashed, "{H|_S}"']
            &
            y(\mathbb A,U,\nu)
                \arrow[d, "H"]
            \\
            \mathcal V
                \arrow[r, hook]
            &
            \mathcal B_{\mathbb B}.
        \end{tikzcd}
        \]
        The dashed map exists if and only if all restrictions of \(H\) along
        arrows of \(S\) lie in \(\mathcal V\). Since
        \(\mathcal B_{\mathbb B}\) is a sheaf, this dashed local membership
        data comes from a section of \(\mathcal V\) exactly when \(H\) itself
        lies in \(\mathcal V\).
    \end{enumerate}
\end{proof}

\begin{definition}[Residual semantic descent]
\label{def:residual-semantic-descent}
    Let
    \[
        \mathcal C=(\mathcal C_{\mathbb A})_{\mathbb A\in\mathsf{Inp}}
    \]
    be a replete global family of full subcategories
    \[
        \mathcal C_{\mathbb A}
        \subseteq
        \mathsf{Mly}^{\mathrm{str}}_{\mathbb A,\mathbb B}.
    \]
    Suppose the generated local family
    \[
        V_{\mathcal C,\mathbb A}
        =
        \bigvee_{P\in\mathcal C_{\mathbb A}}
        \operatorname{Beh}(P)
        \hookrightarrow
        \mathbf{Beh}^{\mathrm{can}}_{\mathbb A,\mathbb B}
    \]
    has been formed for each \(\mathbb A\), using a small representative set of
    distinct semantic-image subobjects.

    We say that \(\mathcal C\) satisfies \textbf{residual semantic descent} if
    the generated family
    \[
        V_{\mathcal C}
        =
        (V_{\mathcal C,\mathbb A})_{\mathbb A\in\mathsf{Inp}}
    \]
    is a sheaf-local global coequational theory.
\end{definition}

Equivalently, residual semantic descent says that membership in the
coequational theory generated by \(\mathcal C\) can be checked after
restriction along canonical residual behaviour covers. This is not an
additional algebraic closure operation on machines; it is the descent condition
saying that the generated residual coequational theory is already a subsheaf
for the canonical residual behaviour topology.

\subsection{Sheaf-local Birkhoff theorem}
\label{subsec:sheaf-local-birkhoff-theorem}

\begin{theorem}[Sheaf-local global coequational Birkhoff theorem]
\label{thm:sheaf-local-global-coequational-birkhoff}
    Let
    \[
        \mathcal C=(\mathcal C_{\mathbb A})_{\mathbb A\in\mathsf{Inp}}
    \]
    be a replete global family of full subcategories
    \[
        \mathcal C_{\mathbb A}
        \subseteq
        \mathsf{Mly}^{\mathrm{str}}_{\mathbb A,\mathbb B}.
    \]
    Then \(\mathcal C\) is of the form
    \[
        \mathcal C_{\mathbb A}
        =
        \mathsf{Mod}^{g}(V)_{\mathbb A}
    \]
    for a sheaf-local global coequational theory
    \[
        V=(V_{\mathbb A})_{\mathbb A\in\mathsf{Inp}}
    \]
    if and only if the following conditions hold:
    \begin{enumerate}
        \item each \(\mathcal C_{\mathbb A}\) is closed under strict
        homomorphic preimages;
        \item each \(\mathcal C_{\mathbb A}\) is closed under
        regular-epimorphic quotients;
        \item each \(\mathcal C_{\mathbb A}\) is closed under small
        coproducts;
        \item the global family \(\mathcal C\) is closed under input pullback
        along functors in \(\mathsf{Inp}\);
        \item \(\mathcal C\) satisfies residual semantic descent.
    \end{enumerate}
    When these conditions hold, the defining sheaf-local theory is
    \[
        V_{\mathcal C,\mathbb A}
        =
        \bigvee_{P\in\mathcal C_{\mathbb A}}
        \operatorname{Beh}(P).
    \]
\end{theorem}

\begin{proof}
    \leavevmode
    \begin{enumerate}
        \item \textbf{Prove necessity of the ordinary closure conditions.}
        Suppose
        \[
            \mathcal C=\mathsf{Mod}^{g}(V)
        \]
        for a sheaf-local global coequational theory
        \[
            V=(V_{\mathbb A})_{\mathbb A\in\mathsf{Inp}}.
        \]
        Conditions \((1)\)--\((4)\) are exactly the closure properties proved
        in Theorem~\ref{thm:global-coequational-closure}.

        \item \textbf{Identify the generated theory.}
        For every
        \[
            P\in\mathcal C_{\mathbb A},
        \]
        the semantic map of \(P\) factors through \(V_{\mathbb A}\). Hence
        \[
            \operatorname{Beh}(P)\leq V_{\mathbb A}.
        \]
        Taking the join over all such \(P\), one obtains
        \[
            V_{\mathcal C,\mathbb A}\leq V_{\mathbb A}.
        \]

        Conversely,
        \[
            V_{\mathbb A}
        \]
        is itself a model of \(V_{\mathbb A}\), regarded as the restricted
        canonical object. By
        Lemma~\ref{lem:semantic-map-of-a-restricted-canonical-subobject},
        \[
            \operatorname{Beh}(V_{\mathbb A})=V_{\mathbb A}.
        \]
        Since
        \[
            V_{\mathbb A}\in\mathcal C_{\mathbb A},
        \]
        it follows that
        \[
            V_{\mathbb A}\leq V_{\mathcal C,\mathbb A}.
        \]
        Therefore
        \[
            V_{\mathcal C,\mathbb A}=V_{\mathbb A}
        \]
        for every \(\mathbb A\).

        The comparison is
        \[
        \begin{tikzcd}[column sep=large, row sep=large]
            V_{\mathbb A}
                \arrow[r, equal]
                \arrow[dr, hook, "j_{V_{\mathbb A}}"']
            &
            \operatorname{Beh}(V_{\mathbb A})
                \arrow[d, hook]
            \\
            &
            \mathbf{Beh}^{\mathrm{can}}_{\mathbb A,\mathbb B}.
        \end{tikzcd}
        \]

        \item \textbf{Prove residual semantic descent.}
        Since
        \[
            V_{\mathcal C}=V
        \]
        and \(V\) is sheaf-local, the generated family
        \[
            V_{\mathcal C}
        \]
        is a sheaf-local global coequational theory. Hence \(\mathcal C\)
        satisfies residual semantic descent.

        \item \textbf{Prove sufficiency.}
        Conversely, assume conditions \((1)\)--\((5)\). By
        Theorem~\ref{thm:global-coequational-birkhoff}, conditions
        \((1)\)--\((4)\) imply that the generated family
        \[
            V_{\mathcal C}
        \]
        is a global coequational theory and that
        \[
            \mathcal C
            =
            \mathsf{Mod}^{g}(V_{\mathcal C}).
        \]
        By condition \((5)\), \(V_{\mathcal C}\) is sheaf-local. Therefore
        \[
            \mathcal C
        \]
        is definable by a sheaf-local global coequational theory.

        \item \textbf{Identify the defining theory.}
        The defining theory is the generated theory
        \[
            V_{\mathcal C,\mathbb A}
            =
            \bigvee_{P\in\mathcal C_{\mathbb A}}
            \operatorname{Beh}(P),
        \]
        by construction.
    \end{enumerate}
\end{proof}

\begin{corollary}[Structured sheaf-local Birkhoff theorem]
\label{cor:structured-sheaf-local-birkhoff}
    Suppose \(\mathbb B\) is a \(\mathbb T\)-structured output category. A
    replete global family
    \[
        \mathcal C=(\mathcal C_{\mathbb A})_{\mathbb A\in\mathsf{Inp}}
    \]
    is of the form
    \[
        \mathcal C_{\mathbb A}
        =
        \mathsf{Mod}^{g}(V)_{\mathbb A}
    \]
    for a sheaf-local global \(\mathbb T\)-coequational theory \(V\) if and
    only if:
    \begin{enumerate}
        \item each \(\mathcal C_{\mathbb A}\) is closed under strict
        homomorphic preimages;
        \item each \(\mathcal C_{\mathbb A}\) is closed under
        regular-epimorphic quotients;
        \item each \(\mathcal C_{\mathbb A}\) is closed under small
        coproducts;
        \item the global family is closed under input pullback along functors
        in \(\mathsf{Inp}\);
        \item the global family is closed under the operations induced by
        \(\mathbb T\), pointwise at each input category \(\mathbb A\);
        \item the generated global \(\mathbb T\)-coequational theory satisfies
        residual semantic descent.
    \end{enumerate}
\end{corollary}

\begin{proof}
    \leavevmode
    \begin{enumerate}
        \item \textbf{Use the structured global Birkhoff theorem.}
        Conditions \((1)\)--\((5)\), without the descent condition, are
        precisely the closure conditions of the structured global Birkhoff
        theorem. Hence they imply that the generated global theory
        \[
            V_{\mathcal C}
        \]
        is a global \(\mathbb T\)-coequational theory and that
        \[
            \mathcal C=\mathsf{Mod}^{g}(V_{\mathcal C}).
        \]

        \item \textbf{Add sheaf-locality.}
        Condition \((6)\) says exactly that
        \[
            V_{\mathcal C}
        \]
        is sheaf-local. Therefore \(\mathcal C\) is definable by a
        sheaf-local global \(\mathbb T\)-coequational theory.

        \item \textbf{Prove necessity.}
        Conversely, if
        \[
            \mathcal C=\mathsf{Mod}^{g}(V)
        \]
        for a sheaf-local global \(\mathbb T\)-coequational theory \(V\), then
        the ordinary structured closure conditions follow from
        Theorem~\ref{thm:global-coequational-closure} and
        Theorem~\ref{thm:T-coequational-closure}, applied pointwise. The same
        argument as in
        Theorem~\ref{thm:sheaf-local-global-coequational-birkhoff} gives
        \[
            V_{\mathcal C}=V,
        \]
        so the generated theory satisfies residual semantic descent.
    \end{enumerate}
\end{proof}

\begin{remark}[Presheaf-level and sheaf-level globality]
\label{rem:presheaf-level-and-sheaf-level-globality}
    The global coequational Birkhoff theorem is a presheaf-level statement over
    the residual context category. Its globality condition is stability under
    residual-context restriction.

    The sheaf-local theorem adds a descent condition. A sheaf-local theory is
    not merely closed under restriction; it also satisfies the converse
    locality principle that membership in the theory is detected after
    restriction along canonical residual behaviour covers.
\end{remark}

\begin{remark}[Recovery of the pointwise presentation]
\label{rem:recovery-of-pointwise-presentation}
    When \(\mathcal E=\mathbf{Set}\), or in a well-pointed situation where
    global elements are sufficient for the intended applications, the
    generalized context
    \[
        (\mathbb A,U,\nu)
    \]
    may be specialized to
    \[
        U=1,
        \qquad
        \nu:1\to A_0.
    \]
    Then one recovers the simpler pointed notation
    \[
        (\mathbb A,v),
    \]
    and a residual context morphism becomes a pair
    \[
        (F,r):(\mathbb A',v')\to(\mathbb A,v),
        \qquad
        r:F_0(v')\to v.
    \]
    The generalized formulation above is the version that remains correct in
    arbitrary Grothendieck toposes.
\end{remark}

